\section{Protocol: audited perturb-and-validate}
\label{sec:protocol}

% source: PAPER_SKELETON Sec 2; EVIDENCE_protocol_prereg P1-P8.

\paragraph{We plant a statement into the model's own trace and formally
validate what it does next.} The unit of analysis is a single statement
inserted into a model's in-stream reasoning trace on a formally checkable
first-order-logic task in the PrOntoQA family \citep{saparov2023prontoqa,
saparov2023ood}, after which the model continues with no external
correction. Two audits bracket every row. Before generation, each planted
statement is truth-audited against rule closure in the closed world, so we
know whether the plant is actually true or false rather than assuming it.
After generation, a formal validator assigns each continuation to one of a
set of mutually exclusive, precedence-ordered outcome classes. The classes,
in precedence order, are generation-failed, unparsed, derailed (the final
answer differs from the target), injection-dependent (at least one step is
derivable only via the planted atom, but the final answer still matches the
target), parroted (the final answer matches despite an underivable step),
and closure-valid completion (every step is entailed by the true premises
and prior valid steps under rule closure).
% source: EVIDENCE_protocol_prereg P1 (validatorPlan L115-117)
Two facts anchor these definitions for the rest of the paper.
Closure-valid is a closure-based standard, so a single continuation
sentence may legally jump several hops; it is not a strict replay. And the
word ``recovery'' is reserved throughout for one defined metric, strict
repair, introduced below; we otherwise say the model rejected, absorbed,
flagged, or reused a plant.
% source: EVIDENCE_protocol_prereg P1 caveat (validatorPlan L45, L123-125); STYLE_SHEET Sec 3

\paragraph{The truth audit is load-bearing, not decorative.} An earlier
version of this instrument treated a ``wrong-category'' perturbation family
as a family of falsehoods. Under the audit, 355 of its 390 members are
actually true and only 35 are false, so a headline built on it would have
measured behavior on mostly-true statements.
% source: EVIDENCE_protocol_prereg P2 (STRICT_METRICS_REPORT L257, L328)
On the matched-distance grid used here the audit is clean: every
planted-false cell is audited false and every anchor cell audited true.
% source: EVIDENCE_protocol_prereg P2 (EXPD_FULL_REPORT integrity L182-195)
We report the 355-of-390 inversion in the protocol, not as a result about
models, because it is the reason the audit exists.

\paragraph{The design matrix crosses the factors so they are pinned, not
confounded.} The matched-distance grid crosses truth status, polarity,
derivational usability, and measured inferential distance across 24 cells,
with parse gaps at most 0.67 percentage points across matched cells, so
visibility, usability, and polarity can be read separately rather than
being tangled together (Table~\ref{tab:design}).
% source: PAPER_SKELETON 2.2; EVIDENCE_protocol_prereg P6 (EXPD grid)
This is the visual answer to the reviewer question of whether the effects
are confounded: each row of Table~\ref{tab:design} holds the other factors
fixed.

\paragraph{Distance is measured per row, not assumed from the design
label.} We do not trust the designer's intended distance. A mechanical
breadth-first search over rule applications computes, for each row, the
shortest number of hops from the visible-prefix state to the evidence that
would refute the plant. On the generated worlds used for the controlled
results, designed and measured distance agree and the global design point
is clean, with zero of 450 rows carrying a finite refutation distance.
% source: EVIDENCE_protocol_prereg P3 (RETRO_DISTANCE L88-93)
The audit also exposes why a legacy distance ladder cannot support a
dose-response reading: its designed one-hop condition measures at distance
zero on 14.9 percent of rows and at distance two or more on 8.4 percent,
and its multi-hop rungs are heavily smeared, so any gradient over the
designed label there is an artifact and we re-bin by measured distance.
% source: EVIDENCE_protocol_prereg P3 (RETRO 76.7/14.9/8.4, neghop smearing)
Measured distance is rule-BFS distance in the code-native fragment, not hop
count in the model's own written proof; we keep that distinction explicit.

\paragraph{One box fixes every dependent-variable definition and the
backend rule.} To keep the instrument honest we fix, in one place, the
definitions of the stated-complement rejection measure, closure-valid and
hop-sound validity, the strict-repair rate, and derivational
injection-dependence, together with two house rules that bind every number
in the paper.
% source: PAPER_SKELETON 2.4; EVIDENCE_protocol_prereg P8
The first rule is that generation never mixes inference backends within a
comparison, because the verbalized-doubt measure is backend-sensitive by 8
to 15 percentage points; the second is that doubt is reported only as a
contrast, never as an absolute rate, for the same reason.
% source: EVIDENCE_protocol_prereg P8 (EXPE L9 DV registry; backend sensitivity)
The verbalized-doubt measure is a lexical proxy that we demoted to
descriptive before review, and Section~\ref{sec:instrument} shows why.

\paragraph{Planting is justified because the modal natural error lives
exactly where in-stream checking is hardest.} A design that plants errors
is only interesting if it plants the errors models actually make. We
audited every unperturbed rollout in a corpus of 3{,}715 gold rollouts and
found 236 natural invalid entity-facts.
% source: EVIDENCE_protocol_prereg P7 (NATURAL_ERRORS n=236, 3715 rollouts)
The planted families cover 93 percent of the natural invalid entity-facts,
and the single modal natural error, at 85 percent, is an
affirmative-unentailed category whose refutation is more than eight hops
away, which is exactly the global design point; locally refutable natural
errors are rare, at 8 percent.
% source: EVIDENCE_protocol_prereg P7 (93% coverage, 85% modal cell, 8% local)
The natural distribution also reproduces the paper's central asymmetry,
with naturally usable errors consumed downstream about eleven times more
often than inert ones (Section~\ref{sec:usability}).
% source: EVIDENCE_protocol_prereg P7 (0.54 usable vs 0.047 inert natural)
One honest gap travels with this: the single most frequent natural error
\emph{object} is a fabricated or mis-spliced rule rather than a false
entity-fact, which no planted family emulates except in inert form, so we
scope the paper's claims to entity-fact falsehoods.
% source: EVIDENCE_protocol_prereg P7 caveat 1 (1118 unentailed rules, 41% load-bearing)

\paragraph{The grid is auditable end-to-end and fails closed.} The
matched-distance grid generated 2{,}880 worlds across 360 families with zero
world-audit failures and 11{,}577 validated rows, all with unique run
identifiers. Injection audits reject fail-closed rather than silently
passing a mis-audited plant, including 57 rows dropped for a
measured-distance mismatch.
% source: EVIDENCE_protocol_prereg P6 (funnel L159-195); PAPER_SKELETON 2.6
The multi-model replication pipeline re-audits the same way, with zero
world-audit failures on every model (Section~\ref{sec:usability}).
% source: E1_ADJUDICATION Sec 2 (0 world-audit failures all 4 models)

\begin{table}[t]
\centering
\caption{Design matrix (schematic). The matched-distance grid crosses truth
status, polarity, derivational usability, and \emph{measured} refutation
distance, so each factor can be read with the others held fixed. Full
per-cell rates are in Appendix~\ref{app:strict}. \todo{populate the schematic
with the exact 24-cell layout from the matched-distance grid report before submission.}}
\label{tab:design}
\begin{tabular}{@{}llll@{}}
\toprule
Factor & Levels & Held fixed by & Measured? \\
\midrule
Truth status & audited-true / audited-false & closed-world closure audit & yes \\
Polarity & affirmative / negated & matched pairs & yes \\
Usability & usable / inert & derivational need for target & yes \\
Distance $d$ & $0,1,2,3,5,\infty$ & rule-BFS per row & yes (not designed) \\
\bottomrule
\end{tabular}
\end{table}
